% TEMPLATE-MILC-v3.tex - plantilla maestra UNICA de las guias MILC v3.
%
% La usan las tres familias de guias, cada una con su driver:
%   · Grados 6-11  -> scripts/build-guias-g11.py
%   · Bebras       -> scripts/build-guias-bebras.py
%   · Semillero    -> scripts/build-guias-semillero.py
%
% Antes habia tres copias casi identicas (~400 lineas iguales, 6 distintas):
% cada mejora habia que aplicarla tres veces, y en la practica no se hacia
% (el motor de recursos vivia solo en dos de ellas). Lo que varia entre
% familias esta parametrizado en 6 placeholders que cada driver llena
% (se nombran aqui SIN su sintaxis real: el builder reemplaza por texto plano,
% tambien dentro de los comentarios, y un valor multilinea rompe el preambulo):
%
%   HEADER_IZQ        encabezado izquierdo de pagina
%   HEADER_DER        encabezado derecho de pagina
%   PORTADA_ETIQUETA  rotulo sobre el numero grande (GRADO/MOMENTO/MODULO)
%   PORTADA_SERIE     linea de serie de la portada
%   PORTADA_ID        identificador de la guia en la portada
%   PORTADA_CIRCULO   el circulito de grado (°), o vacio si no aplica
%
% El resto de claves las documenta content/guias/_SCHEMA.md. El script valida
% que no queden placeholders sin reemplazar antes de escribir el archivo final.

\documentclass[11pt,letterpaper]{article}
\usepackage[margin=2.0cm]{geometry}
\usepackage[table]{xcolor}
\usepackage{fontspec}
\usepackage[spanish]{babel}
\usepackage{tikz}
% Librerías TikZ para los diagramas de las guías (bloque `recursos:`):
%  · babel  → babel en español vuelve activos caracteres como «>», y TikZ los usa
%             en sus opciones (>=stealth, ->). Sin esta librería, cualquier
%             diagrama con flechas revienta con «\language@active@arg> has an
%             extra }». Debe ir ANTES de que se dibuje nada.
%  · positioning        → sintaxis «below left=2pt and 3pt».
%  · decorations.pathreplacing → llaves ({brace}) para señalar rangos.
\usetikzlibrary{babel,positioning,decorations.pathreplacing}
\usepackage{graphicx}
% Circuitos: el trabajo de electrónica se hace hoy con micro:bit y sus sensores
% integrados (MakeCode, sin cableado), así que no se necesitan esquemáticos.
% Si algún día entran montajes con protoboard, basta descomentar esta línea:
% los diagramas .tex/.tikz declarados en `recursos:` ya se incrustan solos.
% \usepackage{circuitikz}
\usepackage[most]{tcolorbox}
\usepackage{tabularx}
\usepackage{array}
\usepackage{fancyhdr}
\usepackage{titlesec}
\usepackage[document]{ragged2e}  % bandera a la derecha en todo el cuerpo
\usepackage{enumitem}
\usepackage{microtype}

% Helvetica Neue no trae la flecha U+2192, asi que cada «→» escrito literalmente
% en el YAML se compilaba a NADA, en silencio (462 flechas en 61 guias: rutas del
% tipo «paleta → tipografia → lockup» salian rotas). Mapeamos el caracter a su
% version matematica, que si tiene glifo. Asi el autor puede seguir escribiendo
% «→» comodamente en el YAML.
\usepackage{newunicodechar}
\newunicodechar{→}{\ensuremath{\rightarrow}}
% Mismo problema con los simbolos de marca y vinetas: el «✓» de «una palomita ✓»
% salia como cuadrito vacio en el PDF de todas las guias que lo usaban.
\usepackage{pifont}
\newunicodechar{✓}{\ding{51}}
\newunicodechar{✗}{\ding{55}}
\newunicodechar{✔}{\ding{52}}
\newunicodechar{•}{\textbullet}
\newunicodechar{≥}{\ensuremath{\geq}}
\newunicodechar{≤}{\ensuremath{\leq}}

\setmainfont{Helvetica Neue}
\setsansfont{Helvetica Neue}
\setlength{\parindent}{0pt}
\setlength{\parskip}{6pt}
% Sin particion de palabras: en 6.º-7.º una palabra cortada por guion al final
% de linea («Comuni-cacion») frena la lectura mucho mas que una linea corta.
\hyphenpenalty=10000
\exhyphenpenalty=10000
\pretolerance=2000
\tolerance=3000
\setlength{\emergencystretch}{3em}
\linespread{1.10}
\renewcommand{\arraystretch}{1.25}
\setlist{leftmargin=5mm,itemsep=1mm,topsep=1mm}
\newcolumntype{Y}{>{\RaggedRight\arraybackslash}X}

\definecolor{milcMagenta}{HTML}{E6007E}
\definecolor{milcVino}{HTML}{5A0038}
\definecolor{milcTurquesa}{HTML}{0093A5}
\definecolor{milcVerde}{HTML}{58A923}
\definecolor{milcMostaza}{HTML}{E5B400}
\definecolor{milcOcre}{HTML}{B86600}
\definecolor{milcNegro}{HTML}{241816}
\definecolor{milcGris}{HTML}{F7F7F4}
\definecolor{milcLinea}{HTML}{E8E2DD}
\definecolor{milcGradePrimary}{HTML}{0066FF}
\definecolor{milcGradeDark}{HTML}{003D99}
\definecolor{milcGradeSoft}{HTML}{D6E8FF}
\definecolor{milcGradeAccent}{HTML}{CFE7FF}
\definecolor{milcGradeText}{HTML}{FFFFFF}
\color{milcNegro}

\pagestyle{fancy}
\fancyhf{}
\lhead{\small Tecnología e Informática · Grado 6}
\rhead{\small Guía 28 · MILC v3}
\cfoot{\small\thepage}
\renewcommand{\headrulewidth}{0pt}

\newtcolorbox{titlebox}[1]{
  enhanced,colback=#1,colframe=#1,arc=7mm,boxrule=0pt,
  left=7mm,right=7mm,top=3mm,bottom=3mm,
  before skip=8pt,after skip=8pt,
  fontupper=\sffamily\bfseries\Large\color{white}
}

\newtcolorbox{softbox}[3]{
  enhanced,colback=#2,colframe=#1,borderline west={4pt}{0pt}{#1},
  arc=2mm,boxrule=.4pt,left=4mm,right=4mm,top=3mm,bottom=3mm,
  before skip=6pt,after skip=8pt,title={#3},coltitle=white,
  colbacktitle=#1,fonttitle=\sffamily\bfseries,fontupper=\RaggedRight,
  attach boxed title to top left={xshift=3mm,yshift=-2mm},
  boxed title style={arc=2mm, boxrule=0pt}
}

\newtcolorbox{drawbox}[2]{
  enhanced,colback=white,colframe=#1,arc=4mm,boxrule=1pt,
  left=5mm,right=5mm,top=4mm,bottom=4mm,before skip=8pt,after skip=8pt,
  title={#2},coltitle=white,colbacktitle=#1,fonttitle=\sffamily\bfseries,
  fontupper=\RaggedRight,
  attach boxed title to top left={xshift=4mm,yshift=-2mm},
  boxed title style={arc=3mm, boxrule=0pt}
}

\newtcolorbox{infoband}[1]{
  enhanced,colback=#1!8,colframe=#1!50,arc=2mm,boxrule=.5pt,
  left=4mm,right=4mm,top=2mm,bottom=2mm,before skip=4pt,after skip=4pt,
  fontupper=\small\RaggedRight
}

% Puente narrativo entre secciones (contrato editorial Conéctate).
% Da continuidad: "Acabas de X. Ahora vas a Y porque Z."
% Visualmente: banda izquierda en color del grado, texto en itálica pequeña.
\newtcolorbox{puentebox}{
  enhanced,colback=milcGradeSoft,colframe=milcGradeDark,
  borderline west={3pt}{0pt}{milcGradeDark},
  arc=0mm,boxrule=0pt,left=5mm,right=4mm,top=2.5mm,bottom=2.5mm,
  before skip=4pt,after skip=8pt,
  fontupper=\itshape\small\color{milcNegro}\RaggedRight
}
% La flecha va en modo matematico a proposito: Helvetica Neue NO tiene el glifo
% U+2192, asi que \textrightarrow se compilaba a NADA («Missing character: There
% is no → in font Helvetica Neue») y el puente salia sin flecha. En modo
% matematico la flecha viene de la fuente matematica y si se dibuja.
\newcommand{\puente}[1]{%
  \begin{puentebox}%
    \textcolor{milcGradeDark}{$\rightarrow$}\hspace{2mm}#1%
  \end{puentebox}%
}

\newcommand{\linead}[1]{\noindent\color{milcLinea}\rule{#1}{0.5pt}\color{milcNegro}}
\newcommand{\respuesta}[1]{\par\vspace{2mm}\foreach \n in {1,...,#1}{\noindent\color{milcLinea}\rule{\linewidth}{0.5pt}\par\vspace{5mm}}\color{milcNegro}}
\newcommand{\checkbox}{\textcolor{milcGradeDark}{\fbox{\phantom{x}}}\hspace{5pt}}

\newcommand{\phasecard}[4]{
\begin{tcolorbox}[enhanced,colback=#3,colframe=#2,arc=3mm,boxrule=.5pt,height=3.05cm,valign=center,halign=center,left=1.5mm,right=1.5mm,top=2mm,bottom=2mm]
{\sffamily\bfseries\fontsize{22}{24}\selectfont\textcolor{#2}{#1}}\par
{\sffamily\bfseries\small #4}
\end{tcolorbox}
}

% ── Piezas del rediseno 2026 (legibilidad 6.º-11.º) ───────────────────
% Banda de estacion: reemplaza el titlebox macizo. Titulo a la izquierda,
% dato practico (tiempo/modalidad) a la derecha, en una sola linea.
\newcommand{\milcestacion}[2]{%
\begin{tcolorbox}[enhanced,colback=#1,colframe=#1,arc=2mm,boxrule=0pt,
  left=5mm,right=5mm,top=2.6mm,bottom=2.6mm,before skip=11pt,after skip=7pt]
{\sffamily\bfseries\fontsize{15}{18}\selectfont\color{white}#2}
\end{tcolorbox}}

% Tarjeta de paso numerada: sustituye las tablas «Pilar / Aplicacion», que
% eran formato de consulta docente, no de estudiante. El numero va en un
% circulo sobre el borde para que la secuencia se lea de un vistazo.
\newcommand{\milcpaso}[3]{%
\begin{tcolorbox}[enhanced,colback=white,colframe=milcLinea,arc=1.5mm,boxrule=.9pt,
  left=14mm,right=4mm,top=3mm,bottom=3mm,before skip=4pt,after skip=4pt,
  fontupper=\RaggedRight,
  overlay={\node[anchor=north west,circle,fill=#2,text=white,inner sep=0pt,
    minimum size=7.4mm,font=\sffamily\bfseries\small]
    at ([xshift=3.6mm,yshift=-3.2mm]frame.north west) {#1};}]
#3
\end{tcolorbox}}

% Casilla marcable de tamano util con lapiz (la anterior era \fbox{\phantom{x}},
% de alto variable segun el texto que la seguia).
\newcommand{\milcmarca}{\framebox[3.8mm]{\rule{0pt}{3.4mm}}}

% Fila marcable de lista suelta (checks de la estacion 1).
\newcommand{\milccheckrow}[1]{%
\par\vspace{1.8mm}\noindent
\makebox[6.4mm][l]{\raisebox{-.55mm}{\textcolor{milcNegro}{\milcmarca}}}%
\parbox[t]{\dimexpr\linewidth-6.4mm\relax}{\RaggedRight #1}\par}

% Celda marcable dentro de la rubrica (tabla de autoevaluacion). Misma
% sangria francesa, pero pensada para vivir dentro de una columna X.
\newcommand{\milcopcion}[1]{%
\makebox[5.2mm][l]{\raisebox{-.55mm}{\textcolor{milcNegro}{\milcmarca}}}%
\parbox[t]{\dimexpr\linewidth-5.2mm\relax}{\RaggedRight #1}}

% ── Recursos visuales de la guía ──────────────────────────────────────
% Declarados en el bloque `recursos:` del YAML y emitidos por el builder.
% \guiaFigura   → imagen rasterizada o PDF (foto, carta celeste, montaje).
% \guiaDiagrama → fuente TikZ/circuitikz (.tex) incrustada como vector:
%                 es la vía para circuitos y esquemáticos editables.
% #1 = ruta relativa al .tex · #2 = pie (puede ir vacío)
\newcommand{\guiaFigura}[2]{%
\begin{center}
\includegraphics[width=0.88\linewidth,height=0.38\textheight,keepaspectratio]{#1}\\[1.5mm]
{\footnotesize\itshape\textcolor{milcNegro!75}{#2}}
\end{center}
}

\newcommand{\guiaDiagrama}[2]{%
\begin{center}
\input{#1}\\[1.5mm]
{\footnotesize\itshape\textcolor{milcNegro!75}{#2}}
\end{center}
}

\begin{document}
\thispagestyle{empty}

% ── Portada ───────────────────────────────────────────────────────────
% Antes: un rectangulo de color a pagina completa. Se veia bien en pantalla,
% pero en la fotocopiadora del colegio se comia el toner y salia gris sucio.
% Ahora la banda ocupa el tercio superior y el resto va en blanco.
\begin{tikzpicture}[remember picture,overlay]
  \path[fill=milcGradePrimary]
    (current page.north west) rectangle ([yshift=-10.6cm]current page.north east);

  \node[anchor=north west,text=milcGradeText] at ([xshift=2.0cm,yshift=-1.55cm]current page.north west)
    {\sffamily\bfseries\fontsize{12}{14}\selectfont GRADO};
  \node[anchor=north west,text=milcGradeText,opacity=.85] at ([xshift=2.0cm,yshift=-2.35cm]current page.north west)
    {\sffamily\bfseries\fontsize{8.5}{10}\selectfont SERIE GUÍAS · TECNOLOGÍA E INFORMÁTICA};
  \node[anchor=north east,text=milcGradeText,opacity=.85,align=right] at ([xshift=-2.0cm,yshift=-1.55cm]current page.north east)
    {\sffamily\bfseries\fontsize{9}{12}\selectfont I.E. Sor María Juliana\\Cartago, Valle del Cauca};

  \node[anchor=west,text=milcGradeText] (gradonum) at ([xshift=1.85cm,yshift=-7.1cm]current page.north west)
    {\sffamily\bfseries\fontsize{112}{112}\selectfont 6};
  % El circulito de grado (°) solo aplica a las guias de grado («11°»); en
  % Bebras y Semillero el numero grande es un momento/modulo, no un grado.
  % Se posiciona relativo al nodo `gradonum`, no en coordenadas absolutas.
  \draw[milcGradeText,line width=1.6pt]
    ([xshift=2.0mm,yshift=-4.5mm]gradonum.north east) circle (.15cm);

  \node[anchor=west,text=milcGradeText,text width=11.4cm,align=left] at ([xshift=1.5cm,yshift=0.5cm]gradonum.east)
    {
     {\sffamily\bfseries\fontsize{9}{11}\selectfont\MakeUppercase{Período 3 · Procesador de texto e Internet}}\\[3mm]
     {\sffamily\bfseries\fontsize{21}{25}\selectfont\setlength{\baselineskip}{27pt}Buscar bien\\[4pt]en internet ---\\[4pt]encontrar lo\\[4pt]que necesito\\[4pt]en segundos}\\[3.5mm]
     {\sffamily\bfseries\fontsize{15}{18}\selectfont Guía 28-6-TIC}
    };
\end{tikzpicture}

\vspace*{9.4cm}
\vfill

{\sffamily\bfseries\fontsize{8}{10}\selectfont\textcolor{milcGradeDark}{LO QUE TE LLEVAS HOY}}

\begin{tcolorbox}[enhanced,colback=white,colframe=milcNegro,arc=2mm,boxrule=1.6pt,
  left=5mm,right=5mm,top=4mm,bottom=4mm,before skip=3pt,after skip=12pt,
  fontupper=\RaggedRight\fontsize{11}{15.5}\selectfont]
Un \textbf{laboratorio de búsquedas} en el cuaderno con \textbf{10 búsquedas reales mejoradas}: para cada una anotas la \emph{búsqueda mala} (la primera que se te ocurre, típicamente vaga), \emph{los problemas} que tendría (resultados irrelevantes, demasiados, en otro idioma), \emph{la búsqueda mejorada} (con operadores), y \emph{cuántos resultados realmente útiles} salieron. Más \textbf{una tabla de 6 operadores de búsqueda} con ejemplo de uso.
\end{tcolorbox}

\begin{tcolorbox}[enhanced,colback=milcGris,colframe=milcGris,arc=2mm,boxrule=0pt,
  left=5mm,right=5mm,top=3.5mm,bottom=3.5mm,after skip=12pt]
{\sffamily\bfseries\fontsize{8}{10}\selectfont\textcolor{milcNegro!55}{ESTA GUÍA ES DE}}\\[3mm]
\begin{tabularx}{\linewidth}{X p{3.2cm} p{3.2cm}}
\linead{\linewidth} & \linead{3.0cm} & \linead{3.0cm}\\[-1mm]
{\fontsize{8.5}{10}\selectfont\textcolor{milcNegro!55}{Nombre completo}} &
{\fontsize{8.5}{10}\selectfont\textcolor{milcNegro!55}{Grupo}} &
{\fontsize{8.5}{10}\selectfont\textcolor{milcNegro!55}{Fecha}}\\
\end{tabularx}
\end{tcolorbox}

\vfill

\noindent\textcolor{milcLinea}{\rule{\linewidth}{1pt}}\\[2mm]
{\sffamily\bfseries\fontsize{9.5}{12}\selectfont Autor: Dr. Álvaro Cárdenas Orozco}\\[1mm]
{\sffamily\fontsize{8.5}{11}\selectfont\textcolor{milcNegro!55}{Metodología MILC · apertura ancestral · pensamiento computacional · triángulo Dussel–estoicismo–Floridi}}

\newpage

\section*{Apertura: Buscar bien en internet --- encontrar lo que necesito en segundos}

\begin{softbox}{milcGradeDark}{milcGradeSoft}{Saber ancestral}
\textbf{Antes de Google, cuando doña Mercedes la maestra rural necesitaba información del campo de Cartago, sabía a quién preguntar.} En la vereda La Plata, doña Mercedes no buscaba al azar. Si necesitaba saber sobre \emph{plagas del café}, iba donde \textbf{don Aurelio} que tenía 50 años trabajando en fincas. Si era sobre \emph{tejido del sombrero aguadeño}, iba donde \textbf{doña Carmen} la costurera. Si era sobre \emph{historia del barrio}, donde \textbf{don Lucho} el relojero que vivió ahí 60 años. \textbf{No preguntaba a cualquiera}: preguntaba a la persona correcta. Esa es la clave de buscar bien: \emph{saber a quién preguntar y cómo preguntar}. Hoy preguntamos a Google, Bing o DuckDuckGo, pero la lógica es la misma: \textbf{si preguntas mal, te dan mala respuesta; si preguntas con precisión, encuentras oro en segundos}. Internet tiene miles de millones de páginas. Solo unas pocas tienen lo que necesitas. Saber filtrar es el saber clave del siglo XXI.
\end{softbox}

\begin{softbox}{milcTurquesa}{milcGris}{Saber contemporáneo}
\textbf{Buscar bien en internet} es una habilidad que se aprende. La mayoría de las personas escribe \emph{``palabra suelta''} en Google y se conforma con los primeros resultados. Los que saben buscar usan \textbf{operadores}: símbolos y palabras que le dicen a Google qué excluir, qué incluir, qué tipo de archivo, qué sitio. Con \textbf{6 operadores básicos} puedes pasar de obtener \emph{millones de resultados confusos} a \emph{10 resultados precisos}. \textbf{Los 6 operadores universales:} \texttt{comillas} (para frase exacta), \texttt{site:} (para buscar en un sitio específico), \texttt{filetype:} (para tipo de archivo), \texttt{-palabra} (para excluir), \texttt{OR} (para alternativas), \texttt{rango..} (para fechas o números). Aprenderlos te hace \emph{10 veces más eficiente} en internet. Y los buscadores tienen \textbf{filtros visuales} (tiempo, idioma, tipo) que la mayoría ignora pero que ayudan muchísimo.
\end{softbox}

\begin{softbox}{milcMostaza}{milcGris}{Pregunta-puente}
\textit{Buscas en Google \emph{``perros''}. Sale Wikipedia, criaderos de perros, fotos de perros, historia de perros. Pero tú lo que quieres es \emph{``qué hacer si mi perro tiene fiebre''}. ¿\textbf{Cómo cambiarías la búsqueda} para encontrar exactamente eso?}
\end{softbox}

\begin{softbox}{milcVerde}{milcGris}{Saber y saber hacer}
Al terminar la clase: (1) podrás \textbf{identificar} los 6 operadores de búsqueda más útiles; (2) sabrás \textbf{aplicar} cada uno en búsquedas reales; (3) podrás \textbf{evaluar} si una búsqueda está bien formulada; (4) habrás \textbf{mejorado} 10 búsquedas tuyas con criterio.
\end{softbox}



\small
\begin{tabularx}{\textwidth}{>{\bfseries}p{3.0cm}Y}
\rowcolor{milcGradePrimary!12}Autor & Dr. Álvaro Cárdenas Orozco\\
DBA & Produce contenidos digitales y valida información de fuentes web (MEN, T\&I 6°)\\
Referentes & Saber ancestral de la maestra rural · Lectura crítica · Soberanía informacional\\
Producto final & Un \textbf{laboratorio de búsquedas} en el cuaderno con \textbf{10 búsquedas reales mejoradas}: para cada una anotas la \emph{búsqueda mala} (la primera que se te ocurre, típicamente vaga), \emph{los problemas} que tendría (resultados irrelevantes, demasiados, en otro idioma), \emph{la búsqueda mejorada} (con operadores), y \emph{cuántos resultados realmente útiles} salieron. Más \textbf{una tabla de 6 operadores de búsqueda} con ejemplo de uso.\\
\end{tabularx}
\normalsize

\puente{Hoy haces 4 cosas: identificas el problema de buscar mal, aprendes los 6 operadores, mejoras 10 búsquedas tuyas, y registras todo en el cuaderno.}

\milcestacion{milcGradeDark}{Ruta MILC: así vamos a aprender}
\begin{tabularx}{\textwidth}{>{\centering\arraybackslash}X>{\centering\arraybackslash}X>{\centering\arraybackslash}X>{\centering\arraybackslash}X}
\phasecard{1}{milcVerde}{milcGris}{Escucha\\[-1mm]\small Observo, escucho y pregunto.} &
\phasecard{2}{milcTurquesa}{milcGris}{Entender\\[-1mm]\small Sistematizo y ordeno.} &
\phasecard{3}{milcMagenta}{milcGris}{Hacer\\[-1mm]\small Construyo y aplico.} &
\phasecard{4}{milcVino}{milcGris}{Cierre\\[-1mm]\small Evalúo y reflexiono.}
\end{tabularx}


\puente{Empezamos con 5 búsquedas malas para entender qué falla.}

\milcestacion{milcVerde}{Estación 1 de 3 · Escucha}

\begin{softbox}{milcVerde}{milcGris}{Escena para escuchar}
\textbf{Actividad 1 · ANALIZA --- 5 búsquedas malas y por qué fallan} (12 min · individual). Lee estas 5 búsquedas y \textbf{escribe en el cuaderno qué problema tiene cada una}. \textbf{Búsqueda 1:} ``\texttt{tarea}''. \textbf{Búsqueda 2:} ``\texttt{Colombia}''. \textbf{Búsqueda 3:} ``\texttt{cómo hacer una tarea sobre Bolívar para 6\textsuperscript{o} grado en historia con bibliografía}''. \textbf{Búsqueda 4:} ``\texttt{volcán Galera}'' (con ``e'' en vez de ``a'' en Galeras). \textbf{Búsqueda 5:} ``\texttt{Mexico history}'' (en inglés). \emph{Tip:} las búsquedas pueden ser demasiado vagas, demasiado específicas, con errores ortográficos, en otro idioma, o muy ambiguas.
\end{softbox}

{\sffamily\bfseries\fontsize{8}{10}\selectfont\textcolor{milcNegro!55}{MARCA CUANDO LO HAYAS HECHO}}
\milccheckrow{Leí las 5 búsquedas con calma.}
\milccheckrow{Identifiqué el problema de cada una.}
\milccheckrow{Pensé qué tipo de error es (vaga, ambigua, error, idioma, larga).}

\begin{infoband}{milcVerde}


\textbf{Lo que va al cuaderno (Actividad 1 · ANALIZA).} Título: «Actividad 1 --- 5 búsquedas malas». \textbf{Formato:} lista numerada de 5 búsquedas, cada una con un breve diagnóstico. \textbf{Extensión:} 5 líneas. \textbf{Sabes que terminaste cuando} puedes nombrar el tipo de error de cada una.
\end{infoband}

\puente{Después aprendes los 6 operadores universales con ejemplos.}

\milcestacion{milcTurquesa}{Estación 2 de 3 · Entender}

\begin{softbox}{milcTurquesa}{milcGris}{Lo que vas a entender}
Ahora aprendes los \textbf{6 operadores de oro} que transforman búsquedas malas en búsquedas excelentes. Aprenderlos te toma 15 minutos; \textbf{te ahorra horas} el resto de tu vida.
\end{softbox}

\milcpaso{1}{milcTurquesa}{\textbf{Operador 1 --- Comillas: para buscar frase exacta.} Si escribes \texttt{``volcán Galeras''} entre comillas, Google busca esas 2 palabras \emph{juntas y en ese orden}. Sin comillas, busca cualquier página que tenga las 2 palabras dispersas. \emph{Ejemplo:} buscar \texttt{``parque del bolívar Cartago''} te da el parque de Cartago. Buscar \emph{parque bolívar Cartago} sin comillas te da parques de varios países. \textbf{Operador 2 --- site: para buscar dentro de un sitio.} Si quieres resultados solo de Wikipedia, escribes \texttt{site:es.wikipedia.org volcán Galeras}. Si quieres solo del Ministerio de Educación, \texttt{site:mineducacion.gov.co cómo es el grado 6}. Te filtra por dominio.}
\milcpaso{2}{milcTurquesa}{\textbf{Operador 3 --- filetype: para tipo de archivo.} Si necesitas un PDF, escribes \texttt{filetype:pdf el código colombia infantil}. Te muestra solo PDFs. Útil para tareas que necesitan documentos formales. Otros: \texttt{filetype:doc}, \texttt{filetype:ppt}, \texttt{filetype:xls}. \textbf{Operador 4 --- guion (-) para excluir.} Si buscas \emph{``películas dragones''} pero NO quieres las de Disney, escribes \texttt{películas dragones -Disney}. El guion delante de una palabra la excluye. Útil cuando hay un tema que se mezcla con otro que no te interesa.}
\milcpaso{3}{milcTurquesa}{\textbf{Operador 5 --- OR para alternativas.} Si quieres páginas que mencionen Cartago O Pereira (cualquiera de las 2), escribes \texttt{cultura Cartago OR Pereira}. \emph{OR} en mayúsculas. Útil para comparar o cuando no sabes en cuál de varios lugares está el dato. \textbf{Operador 6 --- rango (..) para fechas o números.} Si buscas eventos entre 1950 y 1960, escribes \texttt{eventos Colombia 1950..1960}. Para precios de un producto entre 100 y 500 mil pesos: \texttt{computador \$100000..\$500000}. Los dos puntos (\texttt{..}) le dicen al buscador que busque un rango.}
\milcpaso{4}{milcTurquesa}{\textbf{Filtros visuales adicionales (los que la gente no usa pero ayudan mucho).} En Google, \emph{después} de buscar, aparecen pestañas: \textbf{Imágenes} (solo imágenes), \textbf{Noticias} (solo notas recientes), \textbf{Videos} (solo videos), \textbf{Mapas}, \textbf{Compras}. Y arriba, \textbf{Herramientas}: te deja filtrar por tiempo (\emph{última semana, último mes, este año}). Estos filtros, combinados con los 6 operadores, te dan superpoderes de búsqueda. \textbf{Pista:} para buscar texto reciente sobre un tema, usa siempre el filtro de tiempo. Los datos viejos pueden estar desactualizados.}

\begin{drawbox}{milcTurquesa}{6 operadores + filtros}
\textbf{Op 1.} ``Frase exacta'' (entre comillas). \\[1mm]
\textbf{Op 2.} site: (dentro de un sitio). \\[1mm]
\textbf{Op 3.} filetype: (tipo de archivo). \\[1mm]
\textbf{Op 4.} -palabra (excluir). \\[1mm]
\textbf{Op 5.} OR (alternativas). \\[1mm]
\textbf{Op 6.} año..año (rango). \\[1mm]
\textbf{Filtros:} Imágenes · Videos · Noticias · Herramientas (tiempo).
\end{drawbox}

\begin{softbox}{milcMostaza}{milcGris}{Errores comunes y su corrección}
\textbf{Errores frecuentes al buscar.} (1) \emph{Buscar 1 sola palabra}: \texttt{historia} te da millones de resultados confusos. Mejor 3-5 palabras concretas. (2) \emph{No revisar ortografía}: \texttt{volcán Galera} (con e) te da resultados extraños. (3) \emph{Quedarse en la primera página}: muchas veces el resultado bueno está en la 2da página. (4) \emph{No usar filtros}: Google ofrece Imágenes, Videos, Tiempo, etc., y mucha gente no los usa. (5) \emph{Confundir publicidad con resultados orgánicos}: los primeros suelen ser anuncios pagados (vienen marcados como \emph{``Anuncio''}). Pasa de largo a los reales.
\end{softbox}

\begin{infoband}{milcTurquesa}


\textbf{Lo que va al cuaderno (Actividad 2 · EXPLICA).} Título: «Actividad 2 --- Los 6 operadores con ejemplo». \textbf{Formato:} tabla de 6 filas, 3 columnas. Columnas: \emph{operador}, \emph{para qué}, \emph{ejemplo real}. \textbf{Extensión:} 6 filas. \textbf{Sabes que terminaste cuando} podrías ayudar a tu mamá a buscar algo en internet en 30 segundos.
\end{infoband}

\puente{Con los operadores claros, mejoras 10 búsquedas tuyas (reales) y las pruebas.}

\milcestacion{milcMagenta}{Estación 3 de 3 · Hacer}

\begin{softbox}{milcMagenta}{milcGris}{Cómo vas a construir tu producto}
\textbf{Actividad 3 · APLICA --- Laboratorio de 10 búsquedas mejoradas} (28 min · individual). Vas a tomar \textbf{10 cosas reales} que tú podrías querer buscar en internet, y para cada una vas a anotar: la búsqueda mala (la primera que se te ocurre) y la búsqueda mejorada con operadores. Si tienes computador disponible, prueba ambas y compara cuántos resultados útiles salen.
\end{softbox}

\milcpaso{1}{milcMagenta}{\textbf{Paso 1 --- Encabezado y tabla.} En una hoja del cuaderno, título: \textbf{``Mi laboratorio de búsquedas''}. Dibuja una tabla de \textbf{10 filas y 4 columnas}. Encabezados: \emph{(a) Lo que quiero buscar} · \emph{(b) Búsqueda mala} · \emph{(c) Búsqueda mejorada (con operador)} · \emph{(d) ¿Funcionó mejor?}.}
\milcpaso{2}{milcMagenta}{\textbf{Paso 2 --- Las 10 búsquedas.} Llena las 10 filas. Para inspiración: \emph{(1) Tarea de historia sobre Bolívar; (2) Cómo se prepara la mazamorra de Cartago; (3) Significado de la palabra ``efímero''; (4) Mejor canal de YouTube para aprender mate; (5) PDF del himno nacional con notas; (6) Foto del volcán Galeras de noche; (7) Documental de la Selección Colombia 1990 (no 2010 ni 2014); (8) Mapa de Cartago con calles; (9) Receta de arepas paisas (NO venezolanas); (10) Resultados de partido Colombia vs Argentina 1993}.}
\milcpaso{3}{milcMagenta}{\textbf{Paso 3 --- Búsqueda mala (col b).} Escribe lo que escribirías si NO supieras de operadores. Ejemplo para (1): \emph{``Bolívar''} o \emph{``tarea Bolívar''}. Es vago: hay miles de Bolívar (la persona, ciudad, departamento, moneda, equipo de fútbol).}
\milcpaso{4}{milcMagenta}{\textbf{Paso 4 --- Búsqueda mejorada (col c).} Aplica al menos UN operador en cada fila. Ejemplo para (1): \texttt{``Simón Bolívar'' biografía site:colombiaaprende.edu.co filetype:pdf}. Para (5): \texttt{``himno nacional Colombia'' notas musicales filetype:pdf}. Para (9): \texttt{arepas paisas -venezolanas receta}. Sé creativo: cada búsqueda usa un operador distinto.}
\milcpaso{5}{milcMagenta}{\textbf{Paso 5 --- Prueba (si hay computador).} Si tienes acceso a un computador en clase, prueba las 10 búsquedas mejoradas en Google. En la columna (d) anota: \emph{``Sí, mejor''} o \emph{``No, igual''} o \emph{``Mejor en imágenes con el filtro Imágenes''}. Si no tienes computador hoy, las pruebas en casa esta semana.}
\milcpaso{6}{milcMagenta}{\textbf{Paso 6 --- Reflexión final.} Debajo de la tabla, escribe en 3 líneas: \emph{``Aprendí que (a)..., (b)..., (c)...''}. Firma con nombre y fecha. La idea es interiorizar que \emph{buscar bien es habilidad}, no don.}

\begin{drawbox}{milcMagenta}{Antes de mostrar tu laboratorio}
\checkbox La tabla tiene 10 filas con sus 4 columnas. \\[1mm]
\checkbox Cada fila usa al menos un operador en la búsqueda mejorada. \\[1mm]
\checkbox Los 6 operadores aparecen al menos una vez en las 10 búsquedas. \\[1mm]
\checkbox Si tuve computador, anoté si funcionó mejor (columna d). \\[1mm]
\checkbox Hay reflexión final de 3 aprendizajes. \\[1mm]
\checkbox Está firmado con nombre y fecha.
\end{drawbox}

\begin{softbox}{milcMagenta}{milcGris}{Plantilla de trabajo}
\textbf{Estructura del laboratorio.} \textit{Tabla 10 × 4:} qué busco, búsqueda mala, búsqueda mejorada (con operador), resultado real. \textit{Debajo:} 3 aprendizajes en lista. \textit{Esquina:} firma y fecha.
\end{softbox}

\begin{infoband}{milcMagenta}


\textbf{Lo que va al cuaderno (Actividad 3 · APLICA).} Título: «Actividad 3 --- Mi laboratorio de 10 búsquedas». \textbf{Formato:} tabla 10x4 + reflexión de 3 aprendizajes. \textbf{Extensión:} 1 página entera. \textbf{Sabes que terminaste cuando} comparando antes y después, ves la diferencia en calidad de búsqueda.
\end{infoband}

\puente{El producto es la tabla de 10 búsquedas mejoradas + la tabla de 6 operadores.}

\milcestacion{milcMostaza}{Producto de la sesión}
\textbf{Laboratorio de 10 búsquedas mejoradas + 6 operadores aplicados}.
\begin{softbox}{milcMostaza}{milcGris}{Criterios de calidad}
(1) La tabla tiene 10 filas completas. (2) Cada búsqueda mejorada usa al menos un operador. (3) Los 6 operadores aparecen al menos 1 vez en el laboratorio. (4) Hay reflexión final con 3 aprendizajes. (5) Está firmado con nombre y fecha.
\end{softbox}

\puente{Antes de salir, comparas resultados: ¿cuánto cambió la calidad después de mejorar la búsqueda?}

\milcestacion{milcVino}{Evaluación liberadora · cinco dimensiones}

\begin{softbox}{milcVino}{milcGris}{Sentido de la evaluación}
Evaluar no es solo revisar una nota. Es reconocer qué aprendí, qué cambió en mí, qué emociones manejé, cómo me relaciono con la ciudad y la comunidad, y qué le dejaré a quien viene detrás.
\end{softbox}

\textbf{Mi reflexión liberadora en cinco dimensiones.} Completa cada frase iniciada:

\small
\begin{tabularx}{\textwidth}{>{\bfseries\RaggedRight\arraybackslash}p{3.4cm}Y>{\RaggedRight\arraybackslash}p{4.6cm}}
\rowcolor{milcVino}\color{white}Dimensión & \color{white}\bfseries Mi respuesta (frame starter) & \color{white}\bfseries Algo que descubrí\\
1. Personal & Lo que más cambió en mí fue \linead{6.5cm} & \linead{4.2cm}\\
2. Emocional & La emoción más fuerte fue \linead{2.5cm} y la regulé \linead{3cm} & \linead{4.2cm}\\
3. Ciudadana & En democracia esto importa porque \linead{6cm} & \linead{4.2cm}\\
4. Local (Cartago) & En mi barrio sería útil para \linead{6.5cm} & \linead{4.2cm}\\
5. Intergeneracional & A mi abuela/abuelo le diría \linead{6.5cm} & \linead{4.2cm}\\
\end{tabularx}
\normalsize

\begin{infoband}{milcVino}
\textbf{Para la próxima sustentación.} Vuelve a esta tabla en seis meses. Verás que las respuestas cambian — no porque lo aprendido cambie, sino porque tú cambias. La evaluación liberadora es una conversación contigo a través del tiempo.
\end{infoband}


\puente{Cerramos con 3 voces que te ayudan a entender por qué buscar bien es saber preguntar bien.}

\milcestacion{milcGradeDark}{Triángulo de pensamiento · cierre filosófico}

\begin{softbox}{milcVino}{milcGris}{Dussel · lente del nosotros}
\textit{Quien sabe preguntar, sabe encontrar. Quien no, depende de lo que le sirvan.}

{\footnotesize\itshape\color{milcNegro!70}--- formulación de esta guía, al modo de Enrique Dussel. No es una cita textual.}

Doña Mercedes preguntaba a la persona correcta porque sabía que \textbf{una buena pregunta encuentra una buena respuesta}. Hoy es lo mismo: \emph{Google es el oráculo más grande del mundo, pero solo responde bien cuando preguntas bien}. Buscar mal te da basura aunque la información buena exista. Buscar bien te abre acceso a un universo. \emph{Saber buscar es saber pensar críticamente sobre la pregunta antes de la respuesta}.

\textbf{¿Estoy entrenado en preguntar bien o solo confío en que la primera respuesta sea la correcta?}
\end{softbox}
\linead{\linewidth}\\[1mm]
\linead{\linewidth}

\begin{softbox}{milcMostaza}{milcGris}{Estoicismo · Marco Aurelio · cuidado interior}
\textit{Lo que cuesta 30 segundos pensar antes, ahorra 30 minutos buscar después.}

{\footnotesize\itshape\color{milcNegro!70}--- formulación de esta guía, al modo de Marco Aurelio. No es una cita textual.}

Mucha gente \emph{pierde el tiempo} buscando en Google: navegando entre 50 resultados malos, leyendo páginas que no responden lo que necesitaban. \textbf{El truco no es velocidad: es pausa}. Pensar 30 segundos qué exactamente quieres y qué operador te puede ayudar. Esa pausa convierte 30 minutos de búsqueda en 30 segundos de hallazgo.

\textbf{¿Hago el ejercicio de pensar antes de buscar, o salto al teclado y espero suerte?}
\end{softbox}
\linead{\linewidth}\\[1mm]
\linead{\linewidth}

\begin{softbox}{milcTurquesa}{milcGris}{Floridi · lente de la infoesfera}
\textit{En la era de internet, la habilidad escasa ya no es el acceso a la información, es saber filtrarla.}

{\footnotesize\itshape\color{milcNegro!70}--- formulación de esta guía, al modo de Luciano Floridi. No es una cita textual.}

Tu abuela, para encontrar información sobre algo, tenía que ir a la \emph{biblioteca municipal} de Cartago, buscar el libro, leer. \emph{Acceso difícil, filtrado fácil} (solo había unos libros). Tú vives al revés: \textbf{acceso fácil, filtrado difícil} (hay miles de millones de páginas). \emph{Tu habilidad clave del siglo XXI es filtrar}. Aprender operadores es la primera dimensión de esa habilidad.

\textbf{¿Estoy desarrollando la habilidad de filtrar, o me ahogo en información sin criterio?}
\end{softbox}
\linead{\linewidth}\\[1mm]
\linead{\linewidth}

\begin{infoband}{milcNegro}
\textbf{Nota para el docente.} Las tres formulaciones son del autor de la guía, no citas textuales de los pensadores: recogen su lente para pensar el tema, no sus palabras. Si vas a citarlos en clase, remite a la obra original.
\end{infoband}

\milcestacion{milcVino}{Mi autoevaluación · marca una casilla por fila}

{\small
\setlength{\extrarowheight}{2.2pt}
\begin{tabularx}{\textwidth}{>{\RaggedRight\arraybackslash\bfseries}p{3.0cm}YYY}
\rowcolor{milcVino}\color{white}\bfseries Criterio & \color{white}\bfseries Lo logré & \color{white}\bfseries En proceso & \color{white}\bfseries Necesito apoyo\\[.6mm]
Comprensión del tema & \milcopcion{Explico las ideas con mis palabras.} & \milcopcion{Reconozco las ideas, falta precisión.} & \milcopcion{Necesito repasar lo básico.}\\[.8mm]
\rowcolor{milcGris}Pensamiento computacional & \milcopcion{Usé los pilares al armar mi producto.} & \milcopcion{Usé algunos pilares.} & \milcopcion{Necesito releer la estación 2.}\\[.8mm]
Aplicación y producto & \milcopcion{Mi producto cumple los criterios.} & \milcopcion{Está armado, con ajustes pendientes.} & \milcopcion{Aún está en construcción.}\\[.8mm]
\rowcolor{milcGris}Reflexión 5 dimensiones & \milcopcion{Respondí cada una con honestidad.} & \milcopcion{Respondí algunas con detalle.} & \milcopcion{Mis respuestas son muy generales.}\\[.8mm]
Triángulo de pensamiento & \milcopcion{Conecté las 3 voces con mi trabajo.} & \milcopcion{Conecté al menos una.} & \milcopcion{No las relaciono con mi proceso.}\\[.8mm]
\end{tabularx}}

\vspace{3mm}
\textbf{Hoy entendí que:}\\[1mm]
\linead{\linewidth}\\[1mm]
\linead{\linewidth}

\textbf{Mi compromiso para la próxima sesión es:}\\[1mm]
\linead{\linewidth}\\[1mm]
\linead{\linewidth}

\begin{softbox}{milcMostaza}{milcGris}{Cita de despedida · Marco Aurelio}
\textit{``El obstáculo en el camino se vuelve el camino. Lo que estorba al camino se vuelve camino.''}\\[1mm]
Cada error de hoy, bien observado, se convierte en pieza del camino que pisarás mañana.
\end{softbox}

\small
\begin{tabularx}{\textwidth}{>{\bfseries}p{3.6cm}Y}
\rowcolor{milcGradeDark!12}Nombre completo: & \linead{10cm}\\
Fecha: & \linead{4cm} \quad Curso/grupo: \linead{4cm}\\
Mi firma: & \linead{9cm}\\
\end{tabularx}
\normalsize

\begin{infoband}{milcGradeDark}
\textbf{Para la próxima sesión.} Esta semana, cada vez que busques algo en Google, usa \emph{al menos un operador}. Tu cerebro lo va a interiorizar rápido. La próxima clase aprendes a \textbf{evaluar fuentes}: ¿cómo sabes si una página dice la verdad o miente? Es la otra cara de buscar bien: saber distinguir el oro del bronce. Hoy aprendiste a buscar; la próxima aprenderás a verificar.
\end{infoband}

\end{document}
